\documentclass[12pt,a4paper]{article}

\usepackage[utf8]{inputenc}
\usepackage[T1]{fontenc}
\usepackage[brazil]{babel}
\usepackage{amsmath}
\usepackage{amssymb}
\usepackage{geometry}
\usepackage{enumitem}

\geometry{margin=2.5cm}
\setlength{\parskip}{0.7em}
\setlength{\parindent}{0pt}

\title{Cálculo e Fluxo de Criação de Tarefas}
\author{Antikaos}
\date{\today}

\begin{document}

\maketitle

\section{Ideia geral}

Cada deadline pode gerar tarefas planejadas. A geração respeita:

\begin{itemize}
    \item o tipo da deadline: \textit{contínua} ou \textit{pontual};
    \item a data final da deadline;
    \item a opção \textit{Começar a me dedicar};
    \item a recorrência escolhida, no caso das deadlines contínuas;
    \item a disponibilidade diária do usuário;
    \item a prioridade dada por proximidade da entrega e dificuldade percebida.
\end{itemize}

Nenhuma tarefa é gerada no próprio dia da entrega da deadline. Nesse dia, a deadline aparece no calendário, mas não cria uma tarefa de dedicação.

\section{Início da geração}

Para cada deadline $d_i$, define-se uma data inicial de geração $G_i$:

\[
G_i =
\begin{cases}
t, & \text{se o modo for ``hoje''} \\
S_i, & \text{se o modo for ``amanhã''} \\
\max(S_i, \operatorname{dia}(t)), & \text{se o modo for ``data específica''} \\
t, & \text{se o modo for ``após deadline'' e a dependência estiver concluída} \\
\varnothing, & \text{se o modo for ``após deadline'' e a dependência não estiver concluída}
\end{cases}
\]

Se $G_i = \varnothing$, nenhuma tarefa é criada para a deadline naquele momento.

\subsection*{Legenda}

\begin{description}[leftmargin=2.2cm]
    \item[$G_i$] Data inicial efetiva de geração das tarefas da deadline $i$.
    \item[$t$] Momento atual.
    \item[$S_i$] Data salva para início de dedicação, quando aplicável.
    \item[$\operatorname{dia}(t)$] Início do dia de $t$.
\end{description}

\section{Intervalo válido}

As tarefas só podem ser geradas em dias anteriores ao dia da entrega:

\[
\operatorname{dia}(G_i) \leq x < \operatorname{dia}(T_i)
\]

\subsection*{Legenda}

\begin{description}[leftmargin=2.2cm]
    \item[$x$] Dia candidato para criação de tarefa.
    \item[$T_i$] Data e hora final da deadline $i$.
\end{description}

\section{Deadlines contínuas}

Para uma deadline contínua, o app percorre os dias do intervalo válido e aplica a regra de recorrência.

\subsection{Recorrência diária}

Se a recorrência for ``todos os dias'', todo dia válido gera uma tarefa:

\[
R_i(x) = 1
\]

\subsection{Recorrência a cada dois dias}

Se a recorrência for ``a cada dois dias'', gera tarefa quando a distância para o início é par:

\[
R_i(x) =
\begin{cases}
1, & \text{se } \delta(x, G_i) \equiv 0 \pmod{2} \\
0, & \text{caso contrário}
\end{cases}
\]

\subsection{Recorrência por dias da semana}

Se o usuário escolhe dias da semana, uma tarefa é gerada quando o dia candidato pertence ao conjunto selecionado:

\[
R_i(x) =
\begin{cases}
1, & \text{se } w(x) \in W_i \\
0, & \text{caso contrário}
\end{cases}
\]

\subsection*{Legenda}

\begin{description}[leftmargin=2.2cm]
    \item[$R_i(x)$] Indicador de geração de tarefa da deadline $i$ no dia $x$.
    \item[$\delta(x, G_i)$] Número de dias entre o início $G_i$ e o dia $x$.
    \item[$w(x)$] Dia da semana de $x$.
    \item[$W_i$] Conjunto de dias da semana escolhidos pelo usuário.
\end{description}

\section{Deadlines pontuais}

Para uma deadline pontual, o usuário informa:

\begin{itemize}
    \item número de sessões planejadas;
    \item tempo de dedicação por sessão.
\end{itemize}

Seja $H_i$ o conjunto de dias válidos entre $G_i$ e o dia anterior à entrega:

\[
H_i = \{x \mid \operatorname{dia}(G_i) \leq x < \operatorname{dia}(T_i)\}
\]

O número efetivo de sessões é:

\[
u_i = \min(P_i, |H_i|)
\]

As $u_i$ sessões são distribuídas uniformemente ao longo de $H_i$. Para cada sessão $j$, com $0 \leq j < u_i$, o índice escolhido é:

\[
\ell_j =
\operatorname{round}
\left(
\frac{j(|H_i|-1)}{\max(u_i-1,1)}
\right)
\]

\subsection*{Legenda}

\begin{description}[leftmargin=2.2cm]
    \item[$H_i$] Conjunto de dias disponíveis antes da entrega.
    \item[$P_i$] Número de sessões planejadas pelo usuário.
    \item[$u_i$] Número efetivo de sessões que cabem no intervalo.
    \item[$\ell_j$] Índice do dia escolhido para a sessão $j$ dentro de $H_i$.
\end{description}

\section{Reserva discreta no dia da entrega}

No dia final de uma deadline, nenhuma tarefa de dedicação é criada. Porém, para o cálculo de disponibilidade do dia, cada deadline que termina naquela data reserva discretamente $60$ minutos:

\[
B(x) = 60 \cdot E(x)
\]

\[
E(x) =
\left|
\{d_i \mid \neg C_i \land \operatorname{dia}(T_i) = x\}
\right|
\]

\subsection*{Legenda}

\begin{description}[leftmargin=2.2cm]
    \item[$B(x)$] Minutos reservados discretamente no dia $x$ por entregas finais.
    \item[$E(x)$] Número de deadlines não concluídas que terminam no dia $x$.
\end{description}

\section{Disponibilidade diária}

Para cada dia $x$, o usuário possui um tempo disponível $M(x)$, em minutos. Primeiro, subtraem-se a reserva discreta e as tarefas pontuais, que têm preferência:

\[
L(x) =
\max
\left(
M(x) - B(x) - F(x),
0
\right)
\]

\[
F(x) =
\sum_{g \in P(x)} m_g
\]

\subsection*{Legenda}

\begin{description}[leftmargin=2.2cm]
    \item[$M(x)$] Minutos disponíveis do usuário no dia $x$.
    \item[$L(x)$] Minutos restantes para distribuir entre tarefas contínuas.
    \item[$P(x)$] Conjunto de tarefas pontuais candidatas no dia $x$.
    \item[$m_g$] Minutos fixos da tarefa pontual $g$, informados pelo usuário.
\end{description}

\section{Peso de prioridade para tarefas contínuas}

Cada deadline contínua candidata recebe um peso $w_i$:

\[
w_i =
\begin{cases}
4, & \text{se está próxima e é difícil} \\
3, & \text{se está próxima e não é difícil} \\
2, & \text{se não está próxima e é difícil} \\
1, & \text{se não está próxima e não é difícil}
\end{cases}
\]

Uma deadline é considerada próxima quando faltam até $3$ dias para sua entrega:

\[
\Delta_i(x) \leq 3
\]

Ela é considerada difícil quando:

\[
q_i \geq 61
\]

O peso usado na distribuição incremental de minutos é:

\[
\omega_i = w_i \cdot \max(q_i, 1)
\]

\subsection*{Legenda}

\begin{description}[leftmargin=2.2cm]
    \item[$w_i$] Peso categórico da deadline contínua $i$.
    \item[$\omega_i$] Peso ponderado por dificuldade.
    \item[$q_i$] Dificuldade percebida da deadline $i$.
    \item[$\Delta_i(x)$] Dias restantes entre o dia $x$ e a entrega da deadline $i$.
\end{description}

\section{Distribuição de minutos para tarefas contínuas}

Cada tarefa contínua criada precisa ter pelo menos $20$ minutos e no máximo $90$ minutos. Os minutos são sempre múltiplos de $5$.

Seja $C(x)$ o conjunto de tarefas contínuas candidatas no dia $x$. Primeiro, se $L(x) < 20$, nenhuma tarefa contínua é criada:

\[
L(x) < 20 \Rightarrow C_{\text{criada}}(x) = \varnothing
\]

Caso haja tempo, o número máximo de tarefas contínuas criáveis é:

\[
K(x) =
\left\lfloor
\frac{L(x)}{20}
\right\rfloor
\]

As candidatas são ordenadas por:

\begin{enumerate}
    \item maior peso $w_i$;
    \item maior dificuldade percebida $q_i$;
    \item data de entrega mais próxima.
\end{enumerate}

As primeiras $K(x)$ candidatas recebem $20$ minutos iniciais:

\[
m_i = 20
\]

O restante é:

\[
L_r(x) = L(x) - 20 \cdot |C_{\text{criada}}(x)|
\]

Enquanto $L_r(x) \geq 5$, o app adiciona blocos de $5$ minutos à tarefa elegível de maior peso ponderado $\omega_i$, respeitando o teto de $90$ minutos:

\[
m_i \leftarrow m_i + 5
\quad \text{se} \quad
m_i + 5 \leq 90
\]

\subsection*{Legenda}

\begin{description}[leftmargin=2.2cm]
    \item[$C(x)$] Conjunto de tarefas contínuas candidatas no dia $x$.
    \item[$C_{\text{criada}}(x)$] Subconjunto de tarefas contínuas efetivamente criadas.
    \item[$K(x)$] Número máximo de tarefas contínuas que cabem no tempo restante.
    \item[$m_i$] Minutos atribuídos à tarefa contínua da deadline $i$.
    \item[$L_r(x)$] Minutos restantes após a alocação mínima.
\end{description}

\section{Criação final da tarefa}

Uma tarefa só é persistida se:

\[
m_i > 0
\]

O título final da tarefa segue a forma:

\[
\text{``Se dedicar } m_i \text{ minutos para } \text{nome da deadline}\text{''}
\]

\section{Preservação de estado}

Quando as tarefas geradas são reconstruídas, o app preserva:

\begin{itemize}
    \item tarefas já concluídas;
    \item identificadores de integração com Apple Lembretes;
    \item atrasos descartados do card de \textit{Tasks atrasadas}.
\end{itemize}

\end{document}
